\subsection{形式幂级数代数中的导子}

设 I 为一个集合，E 为一个结合、交换、含单位元、线性拓扑化、分离且完备的 A-代数，且 \(\mathbf{x} = (x_{i})_{i \in I}\) 为 E 中满足命题 4（IV，第28页）的条件 a) 与 b) 的元素族。设 \(\varphi\) 为由 \(u \mapsto u(x)\) 给出的、把 A{[}{[}I{]}{]} 映入 E 的连续同态；它给 E 赋予一个 A{[}{[}I{]}{]}-模结构。根据 III，第552页，A{[}{[}I{]}{]} 到 A{[}{[}I{]}{]}-模 E 中的一个 A-导子 D 因而是一个 A-线性映射 D: A{[}{[}I{]}{]} → E ，满足关系

\[
\mathrm{D} (u v) = u (\mathbf {x}) \cdot \mathrm{D} (v) + \mathrm{D} (u) \cdot v (\mathbf {x})\tag{15}
\]

对 A{[}{[}I{]}{]} 中的 u, v 成立。

命题 7. --- 设 \((y_{i})_{i\in I}\) 是 E 中元素的一个族。当 I 为无限集时，我们假设 \(y_{i}\) 沿 I 的有限子集的补集所构成的滤子趋于 0 。于是存在唯一的从 A{[}{[}I{]}{]} 到 A{[}{[}I{]}{]}-模 E 中的连续 A-导子 D ，使得 \(\mathrm{D}(\mathbf{X}_{i})=y_{i}\) 对所有 \(i\in I\) 。我们有

\[
\mathrm{D} (u) = \sum_ {i \in I} (\mathrm{D} _ {i} u) (\mathbf {x}) \cdot y _ {i} \quad (u \in \mathrm{A} [ [ I ] ])  .\tag{16}
\]

由于 0 在 E 中有一个由理想组成的基本邻域系，族 \(((\mathrm{D}_i u)(\mathbf{x}) \cdot y_i)_{i \in I}\) 在 E 中可和，对所有 \(u \in A[[I]]\) （《一般拓扑学》，III，第263页，推论 2）。于是公式 (16) 定义了一个 A-线性映射 \(\mathbf{D}: A[[I]] \to E\) 。我们留给读者验证 D 是一个连续导子。设 \(\mathbf{D}_1\) 是从 A{[}{[}I{]}{]} 到 E 中的一个连续 A-导子，使得 \(\mathbf{D}_1(\mathbf{X}_i) = y_i\) 对所有 \(i \in I\) 。连续导子 \(\mathbf{D} - \mathbf{D}_1\) 的核是 A{[}{[}I{]}{]} 的一个包含 1 与诸未定元 \(X_{i}\) 的闭子代数 B。由于多项式代数 \(\mathbf{A}[(\mathbf{X}_{i})_{i\in I}]\) 在 A{[}{[}I{]}{]} 中稠密，我们有 \(B = A[[I]]\) ，从而 \(D_{1} = D\) 。

推论 1. --- 设 \(\Delta\) 为 A-代数 E 的一个连续导子。对每个形式幂级数 \(u \in \mathbf{A}[[\mathrm{I}]]\) ，族 \(((\mathrm{D}_i u)(\mathbf{x}) \cdot \Delta x_i)_{i \in I}\) 可和，并且我们有

\[
\Delta (u (\mathbf {x})) = \sum_ {i \in I} \left(\mathrm{D} _ {i} u\right) (\mathbf {x}) \cdot \Delta x _ {i}.\tag{17}
\]

这由命题 7 得出，因为映射 \(u \mapsto \Delta(u(\mathbf{x}))\) 是从 \(\mathbf{A}[[\mathbf{I}]]\) 到 \(\mathbf{A}[[\mathbf{I}]]\) -模 \(\mathbf{E}\) 中的一个连续导子。

推论 2. --- 导子 \(D_{i}\) 是 A-代数 A{[}{[}I{]}{]} 的唯一连续导子，使得

\[
\mathrm{D} _ {i} (\mathrm{X} _ {i}) = 1, \quad \mathrm{D} _ {i} (\mathrm{X} _ {j}) = 0 \quad f o r \quad j \neq i.\tag{18}
\]

这由推论 1 得出。

推论 3. --- 设 \(f \in \mathbf{A}[[\mathbf{X}_1, \ldots, \mathbf{X}_p]]\) 与 \(g_i \in \mathbf{A}[[\mathbf{Y}_1, \ldots, \mathbf{Y}_q]]\) ，对 \(1 \leqslant i \leqslant p\) 。假设对 \(1 \leqslant i \leqslant p\) ，\(g_i\) 的常数项为零，并令 \(h = f(g_1, \ldots, g_p)\) 。则对 \(1 \leqslant j \leqslant q\) 我们有

\[
\frac {\partial h}{\partial \mathbf {Y} _ {j}} = \sum_ {i = 1} ^ {p} \mathrm{D} _ {i} f (g _ {1}, \dots , g _ {p}) \cdot \frac {\partial g _ {i}}{\partial \mathbf {Y} _ {j}}.\tag{19}
\]

这是推论 1 的特殊情形：\(\mathbf{E} = \mathbf{A}[[\mathbf{Y}_1,\dots ,\mathbf{Y}_q]]\) ，\(x_{i} = q_{i}\) 且 \(\Delta = \partial /\partial \mathbf{Y}_j\) 。

命题 8. --- 设 \(\mathbf{X} = (\mathbf{X}_i)_{i\in I}\) 是由有限个未定元组成的族。

\begin{enumerate}
\def\labelenumi{(\roman{enumi})}
\item
  形式幂级数环 A{[}{[}X{]}{]} 的每个导子都是连续的。
\item
  从多项式环 A{[}X{]} 到形式幂级数环 A{[}{[}X{]}{]} 中的每个导子都以唯一的方式延拓为环 A{[}{[}X{]}{]} 的导子。
\item
  族 \((\mathbf{D}_i)_{i\in I}\) 是 \(\mathbf{A}[[\mathbf{X}]]\) -模的一组基，该模由 \(\mathbf{A}\) -导子（把 \(\mathbf{A}[[\mathbf{X}]]\) 映入其自身）组成。
\end{enumerate}

设 \(\mathfrak{b}_n\) 为所有阶 \(\geqslant n\) 的形式幂级数组成的集合。显然 \(\mathfrak{b}_n\) 是环 \(A[[X]]\) 中的理想，由 \(n\) 次单项式生成。因此 \(\mathfrak{b}_n\) 由 \(n\) 个无常数项的形式幂级数之积的有限和组成；如果 \(D\) 是 \(A[[X]]\) 的一个导子，则我们有

\[
\mathrm{D} (f _ {1} \dots f _ {n}) = \sum_ {i = 1} ^ {n} f _ {1} \dots f _ {i - 1} \mathrm{D} (f _ {i}) f _ {i + 1} \dots f _ {n},
\]

由此立即得出 \(\mathbf{Db}_{n} \subset \mathbf{b}_{n-1}\) 对 \(n \geqslant 1\) 。由于序列 \((\mathfrak{b}_n)_{n \geqslant 0}\) 是 A{[}{[}X{]}{]} 中 0 的一个基本邻域系（IV，第26页注记），D 是连续的，(i) 得证。

设 \(\Delta\) 是从 A{[}X{]} 到 A{[}{[}X{]}{]} 中的一个导子。如前面那样论证，我们可以证明 \(\Delta(h)\) 属于 \(\mathfrak{b}_{n-1}\) ，对每个齐次多项式 \(h\) （次数 \(n \geqslant 1\) ）。现在设 \(u \in \mathbf{A}[[\mathbf{X}]]\) ，并设 \(u_n\) 是其次数为 \(n\) 的齐次分量，取自 \(u\) 。由于 \(\Delta(u_n) \in \mathfrak{b}_{n-1}\) 对 \(n \geqslant 1\) ，族 \((\Delta(u_n))_{n \geqslant 0}\) 在 A{[}{[}X{]}{]} 中可和，从而我们可以定义 A{[}{[}X{]}{]} 到其自身的一个导子 D ，通过

\[
\mathrm{D} (u) = \sum_ {n \geqslant 0} \Delta \left(u _ {n}\right).
\]

我们有 \(D(\mathfrak{b}_n) \subset \mathfrak{b}_{n-1}\) ，因此 \(D\) 是 \(A[[X]]\) 的加法群的一个连续自同态。映射 \(\Phi: (u, v) \mapsto D(uv) - uD(v) - D(u)v\) （从 \(A[[X]] \times A[[X]]\) 到 \(A[[X]]\) ）是连续的，且在 \(A[X] \times A[X]\) 上为零。由于 \(A[X]\) 在 \(A[[X]]\) 中稠密，我们有 \(\Phi = 0\) ；换言之，\(D\) 是 \(A[[X]]\) 到其自身的一个导子，并且是 \(\Delta\) 的延拓。

最后，A{[}X{]} 在 A{[}{[}X{]}{]} 中稠密，且 A{[}{[}X{]}{]} 的每个导子由 (i) 知是连续的；因此存在 \(\Delta\) 到 A{[}{[}X{]}{]} 的导子的唯一延拓。这就证明了 (ii)。

仍需证明 (iii)。公式 (18)（IV，第33页）表明族 \((\mathbf{D}_i)_{i\in I}\) 在 \(\mathbf{A}[[\mathbf{X}]]\) 上线性无关，而公式 (16)（IV，第32页）应用于 \(\mathbf{E} = \mathbf{A}[[\mathbf{X}]]\) 的情形，表明每个 A-导子都是诸 \(\mathbf{D}_i\) 的、系数在 \(\mathbf{A}[[\mathbf{X}]]\) 中的线性组合。

命题 9. --- 设 \((u_{\lambda})_{\lambda \in L}\) 是 A{[}{[}I{]}{]} 中元素的一个无常数项的可和族，D 是 A-代数 A{[}{[}I{]}{]} 的一个连续导子。若 \(f = \prod_{\lambda \in L} (1 + u_{\lambda})\) （IV，第27页，命题 2），则族 \((D u_{\lambda} / (1 + u_{\lambda}))_{\lambda \in L}\) 是

可求和且我们有

\[
\mathrm{D} (f) / f = \sum_ {\lambda \in L} \mathrm{D} (u _ {\lambda}) / (1 + u _ {\lambda}).\tag{20}
\]

如果 \(g\) 与 \(h\) 是 \(\mathbf{A}[[\mathbf{I}]]\) 的两个可逆元素，则

\[
\mathrm{D} (g h) = h. \mathrm{D} g + g. \mathrm{D} h
\]

因此，除以gh后，

\[
\mathrm{D} (g h) / g h = \mathrm{D} (g) / g + \mathrm{D} (h) / h.\tag{21}
\]

对 L 的每个有限子集 M，令 \(f_{M} = \prod_{\lambda \in M} (1 + u_{\lambda})\) 。由 (21)，我们对 Card M 作归纳推出关系

\[
\mathrm{D} \left(f _ {\mathrm{M}}\right) / f _ {\mathrm{M}} = \sum_ {\lambda \in \mathrm{M}} \mathrm{D} \left(u _ {\lambda}\right) / \left(1 + u _ {\lambda}\right).\tag{22}
\]

这证明了当 L 有限时命题 9 成立。现在假设 L 无限，并用 F 表示 L 的有限子集构成的滤过有序集。我们有 \(\lim_{\mathfrak{F}} f_{M} = f\) ，因此（IV，第30页，注记）

\[
\mathrm{D} (f) / f = \lim _ {\mathfrak {F}} \mathrm{D} (f _ {\mathrm{M}}) / f _ {\mathrm{M}}.
\]

于是命题 9 由对 (22) 取极限得出。

